
\documentclass{article}
\usepackage{colortbl}
\usepackage{makecell}
\usepackage{multirow}
\usepackage{supertabular}

\begin{document}

\newcounter{utterance}

\centering \large Interaction Transcript for game `bbh', experiment `navigate', episode 42 with Qwen3.6{-}27B.
\vspace{24pt}

{ \footnotesize  \setcounter{utterance}{1}
\setlength{\tabcolsep}{0pt}
\begin{supertabular}{c@{$\;$}|p{.15\linewidth}@{}p{.15\linewidth}p{.15\linewidth}p{.15\linewidth}p{.15\linewidth}p{.15\linewidth}}
    \# & \multicolumn{2}{c}{Player} && \multicolumn{2}{c}{Game Master} \\
    \hline

    \theutterance \stepcounter{utterance}  
    & & & \multicolumn{4}{p{0.6\linewidth}}{
        \cellcolor[rgb]{0.9,0.9,0.9}{
            \makecell[{{p{\linewidth}}}]{
                \texttt{\tiny{[P1$\langle$GM]}}
                \texttt{Given a series of navigation instructions, determine whether one would end up back at the starting point.} \\
\\ 
\texttt{Q: If you follow these instructions, do you return to the starting point? Turn left. Turn around. Turn left. Take 7 steps. Take 2 steps. Take 4 steps. Take 8 steps.} \\
\texttt{Options:} \\
\texttt{{-} Yes} \\
\texttt{{-} No} \\
\texttt{A: No} \\
\\ 
\texttt{Q: If you follow these instructions, do you return to the starting point? Turn around. Take 1 step. Take 6 steps. Turn around. Take 6 steps. Take 9 steps. Take 1 step.} \\
\texttt{Options:} \\
\texttt{{-} Yes} \\
\texttt{{-} No} \\
\texttt{A: No} \\
\\ 
\texttt{Q: If you follow these instructions, do you return to the starting point? Always face forward. Take 2 steps right. Take 9 steps left. Take 7 steps right.} \\
\texttt{Options:} \\
\texttt{{-} Yes} \\
\texttt{{-} No} \\
\texttt{A: Yes} \\
\\ 
\texttt{Q: If you follow these instructions, do you return to the starting point? Always face forward. Take 7 steps left. Take 2 steps backward. Take 7 steps backward. Take 7 steps backward. Take 3 steps forward.} \\
\texttt{Options:} \\
\texttt{{-} Yes} \\
\texttt{{-} No} \\
\texttt{A:} \\
            }
        }
    }
     \\ \\

    \theutterance \stepcounter{utterance}  
    & \multicolumn{4}{p{0.6\linewidth}}{
        \cellcolor[rgb]{0.561,0.737,0.561}{
            \makecell[{{p{\linewidth}}}]{
                \texttt{\tiny{[P1$\rangle$GM]}}
                \texttt{To determine if the final position is the starting point, we need to track the net displacement in two dimensions: Left/Right (horizontal) and Forward/Backward (vertical). Let's assume the starting position is \$(0, 0)\$.} \\
\\ 
\texttt{1.  **"Always face forward."**} \\
\texttt{    *   This establishes the initial orientation. Let's define "Forward" as the positive Y{-}axis, "Backward" as the negative Y{-}axis, "Right" as the positive X{-}axis, and "Left" as the negative X{-}axis.} \\
\\ 
\texttt{2.  **"Take 7 steps left."**} \\
\texttt{    *   Movement: 7 steps in the negative X direction.} \\
\texttt{    *   Current Position: \$({-}7, 0)\$} \\
\\ 
\texttt{3.  **"Take 2 steps backward."**} \\
\texttt{    *   Movement: 2 steps in the negative Y direction.} \\
\texttt{    *   Current Position: \$({-}7, {-}2)\$} \\
\\ 
\texttt{4.  **"Take 7 steps backward."**} \\
\texttt{    *   Movement: 7 steps in the negative Y direction.} \\
\texttt{    *   Current Position: \$({-}7, {-}2 {-} 7) = ({-}7, {-}9)\$} \\
\\ 
\texttt{5.  **"Take 7 steps backward."**} \\
\texttt{    *   Movement: 7 steps in the negative Y direction.} \\
\texttt{    *   Current Position: \$({-}7, {-}9 {-} 7) = ({-}7, {-}16)\$} \\
\\ 
\texttt{6.} \\
            }
        }
    }
    & & \\ \\

    \theutterance \stepcounter{utterance}  
    & & & \multicolumn{2}{p{0.3\linewidth}}{
        \cellcolor[rgb]{0.9,0.9,0.9}{
            \makecell[{{p{\linewidth}}}]{
                \texttt{\tiny{[GM$|$GM]}}
                \texttt{Target: No} \\
            }
        }
    }
    & & \\ \\

    \theutterance \stepcounter{utterance}  
    & & & \multicolumn{2}{p{0.3\linewidth}}{
        \cellcolor[rgb]{0.9,0.9,0.9}{
            \makecell[{{p{\linewidth}}}]{
                \texttt{\tiny{[GM$|$GM]}}
                \texttt{game\_result = LOSE} \\
            }
        }
    }
    & & \\ \\

\end{supertabular}
}

\end{document}
